\chapter*{A}
\label{chap:a}

\term{Abelian Category}
A category in which every hom-set is an abelian group and composition is bilinear, and in which every morphism has a kernel and a cokernel. In such categories the isomorphism theorems hold and every morphism factors through its image. The categories of abelian groups, of modules over a ring, and of sheaves of abelian groups on a space are all abelian, while general preadditive categories need not be. Abelian categories are the natural setting for homological algebra, where exact sequences, kernels, and cokernels behave as in the module case.

\term{Abelian Group}
A group in which the group operation is commutative, so $ab = ba$ for all elements $a, b$. Abelian groups are the additive groups of rings, modules, and vector spaces.

\term{Abel's Theorem}
A theorem asserting that if a power series converges at a point on its circle of convergence, then its sum has a limit at that point equal to the value of the series there. It also refers to Abel's result that the general quintic equation cannot be solved by radicals.

\term{Abilene Paradox}
A social phenomenon in which a group makes a decision contrary to the private preferences of every member, because each person assumes the others want it. Not a mathematical paradox, but a well-known example of pluralistic ignorance in group decision making.

\term{Absolute Convergence}
A property of a series $\sum a_n$ for which $\sum |a_n|$ converges. An absolutely convergent series is unconditionally convergent, meaning its sum is independent of the order of its terms.

\term{Absolutely Continuous}
A function $f$ on $[a,b]$ is absolutely continuous if for every $\varepsilon > 0$ there is $\delta > 0$ such that the total variation of $f$ over finitely many disjoint subintervals of total length less than $\delta$ is less than $\varepsilon$. By the fundamental theorem of calculus, such functions are exactly those with $f(x) = f(a) + \int_a^x g(t)\,dt$ for a Lebesgue integrable $g$. A measure $\mu$ is absolutely continuous with respect to $\nu$, written $\mu \ll \nu$, if every $\nu$-null set is $\mu$-null; then the Radon--Nikodym theorem yields a density $g$ with $\mu = g\nu$.

\term{Absorbing Element}
An element $z$ of a set with a binary operation such that $z \ast x = x \ast z = z$ for every element $x$. Zero is the absorbing element for multiplication.

\term{Abstract Algebra}
The branch of mathematics studying algebraic structures such as groups, rings, fields, modules, and their homomorphisms, abstracting away from the concrete nature of the elements.

\term{Accumulation Point}
A point $p$ of a topological space such that every neighbourhood of $p$ contains a point of a set $S$ different from $p$. Also called a limit point or cluster point.

\term{Adaptive Mesh Refinement}
A numerical technique that locally refines or coarsens a computational mesh based on error estimates, improving accuracy where needed.

\term{Additive Inverse}
For an element $a$ of a group written additively, the element $-a$ such that $a + (-a) = 0$, the identity element.

\term{Adjoint}
The operator $A^*$ associated with a linear map $A$ between inner product spaces satisfying $\langle Ax, y \rangle = \langle x, A^* y \rangle$ for all vectors $x, y$. Also called the conjugate transpose in the matrix setting.

\term{Adjunction}
A pair of functors $F \dashv G$ with a natural bijection $\operatorname{Hom}(FX, Y) \cong \operatorname{Hom}(X, GY)$, a fundamental concept in category theory.

\term{Albanese}
The Albanese variety of a complex algebraic variety $X$ is the universal abelian variety arising as the target of the Albanese morphism $\mathrm{alb}_X : X \to \mathrm{Alb}(X)$, a covariant analogue of the Jacobian and Picard constructions built from the space of global holomorphic $1$-forms. It is characterized by the universal property that every morphism from $X$ to an abelian variety factors through it. For a smooth projective curve the Albanese variety is canonically isomorphic to its Jacobian, and the Albanese morphism recovers the Abel--Jacobi map.

\term{Algebra}
A vector space equipped with a bilinear multiplication. Also the branch of mathematics studying symbolic manipulation and the structures built from sets and operations.

\term{Algebraic Closure}
A field extension $\bar{k}$ of a field $k$ that is algebraically closed and algebraic over $k$. Every field has an algebraic closure, unique up to isomorphism.

\term{Algebraic Extension}
A field extension $K/F$ in which every element of $K$ is algebraic over $F$, i.e. a root of some non-zero polynomial over $F$.

\term{Algebraic Number}
A complex number that is a root of some non-zero polynomial with rational coefficients. Numbers that are not algebraic are called transcendental.

\term{Algebraic Structure}
A set equipped with one or more operations satisfying a specified collection of axioms, such as a group, ring, field, or lattice.

\term{Algorithm}
A finite, unambiguous procedure that, given an input, produces an output after finitely many elementary steps. Algorithms are the objects of study of computability theory and complexity theory.

\term{Allais Paradox}
An observed choice pattern in which people's decisions between risky gambles violate the independence axiom of expected utility theory, even for careful reasoners. In the classic pair of gambles, most prefer the sure million over a gamble with a small chance of nothing, yet prefer the gamble offering a $0.10$ chance at five million over the one offering a $0.11$ chance at one million, the reversal hinging on the $0.01$ difference in probability. Such behaviour prompted alternative theories of risk such as prospect theory.

\term{Almost Everywhere}
A property that holds on a set whose complement has measure zero. Functions that differ on a null set are considered equal almost everywhere.

\term{Almost Sure Convergence}
The convergence of a sequence of random variables $X_n \to X$ on an event of probability 1, also called convergence with probability one.

\term{Alternating Group}
The group $A_n$ of even permutations of a set of $n$ symbols. It is a normal subgroup of the symmetric group $S_n$ and is simple for $n \geq 5$.

\term{Analysis}
The branch of mathematics built on limits, continuity, differentiation, and integration, including real analysis, complex analysis, and functional analysis.

\term{Analytic Continuation}
The extension of an analytic function to a larger domain by power series or functional relations, unique when the larger domain is connected.

\term{Analytic Function}
A function that is locally representable by a convergent power series. A function is holomorphic if and only if it is analytic.

\term{Annihilator}
For a subspace $S$ of a vector space $V$ with dual $V^*$, the subspace $S^\circ$ of all linear functionals vanishing on $S$.

\term{Annulus}
The region between two concentric circles, $\{z : r < |z| < R\}$ in the complex plane. It is a standard domain for Laurent series expansions.

\term{Antiderivative}
A function $F$ whose derivative is a given function $f$, so that $F' = f$. The set of all antiderivatives of $f$ differs by constants.

\term{Antisymmetric Relation}
A relation $R$ on a set such that $xRy$ and $yRx$ together imply $x = y$. Partial orders are antisymmetric relations.

\term{Archimedean Property}
The property of an ordered field that for any positive elements $a, b$ there exists a natural number $n$ with $na > b$. The real numbers satisfy it; non-archimedean fields do not.

\term{Argument Principle}
The theorem that $\frac{1}{2\pi i}\int_\gamma \frac{f'(z)}{f(z)}\,dz = N - P$, where $N$ is the number of zeros and $P$ the number of poles of a meromorphic function $f$ inside the closed curve $\gamma$.

\term{Arrow's Impossibility Theorem}
A theorem stating that with at least three alternatives, no social welfare function can simultaneously satisfy unanimity, independence of irrelevant alternatives, and non-dictatorship. It shows that no fair, consistent voting method exists for aggregating individual preferences, and it is a central result of welfare economics.

\term{Artinian Ring}
A ring in which every descending chain of left ideals stabilizes, so the descending chain condition holds on left ideals. By Hopkins's theorem a left Artinian ring is left Noetherian and its Jacobson radical is nilpotent, a result with no analogue for general Noetherian rings. Finite rings, matrix rings over division rings, and finite products of these are Artinian, while polynomial rings and most integral domains are not. Artinian rings play a central role in noncommutative ring theory, where the Artinian simple rings are exactly the matrix rings over division rings.

\term{Arzelà--Ascoli Theorem}
A criterion for a family of continuous functions on a compact metric space to be relatively compact in the uniform topology: it suffices that the family be equicontinuous and pointwise bounded. Consequently, a sequence of continuous functions on a compact metric space that is equicontinuous and pointwise bounded has a uniformly convergent subsequence, and $C(X)$ is compact in the uniform topology under these hypotheses.

\term{Asymptote}
A line that a curve approaches arbitrarily closely but does not in general meet, describing the behaviour of the curve at infinity.

\term{Automorphic Form}
A smooth function on a symmetric space, or more generally on the adelic points of a reductive group, that is invariant up to a fixed factor under the left action of an arithmetic or congruence subgroup. Classical modular forms are the case of weight $k$ on the upper half-plane, transforming as $f\bigl(\frac{az+b}{cz+d}\bigr) = (cz+d)^k f(z)$. Automorphic forms generate automorphic representations by right translation, and their Hecke eigenvalues are tied to $L$-functions, making them central objects in the Langlands program.

\term{Automorphism}
An isomorphism from a mathematical structure to itself. The set of automorphisms of a structure forms a group under composition.

\term{Axiom}
A statement taken to be self-evidently true, serving as a starting point for further reasoning. Modern mathematics relies on explicit axiom systems such as the Zermelo--Fraenkel axioms.

\term{Axiom of Choice}
The set-theoretic axiom stating that for any collection of non-empty sets there is a choice function selecting one element from each. It is independent of the Zermelo--Fraenkel axioms.

